\documentclass[a4paper]{article}

% Basic packages
% Español (México) con XeLaTeX: fontspec para fuentes Unicode, polyglossia
% para la división silábica y los nombres del documento en la variante mexicana.
\usepackage{fontspec}
\usepackage{polyglossia}
\setdefaultlanguage[variant=mexican]{spanish}
% Los nombres chinos van como «pinyin (original)»: xeCJK da a los caracteres
% Han su propia fuente; sin ella XeLaTeX los omite con un aviso.
% FandolSong no cubre algunos caracteres de nombres (U+73FA); WenQuanYi Zen Hei sí.
\usepackage[AutoFallBack=true]{xeCJK}
\setCJKmainfont{FandolSong-Regular.otf}
\setCJKfallbackfamilyfont{\CJKrmdefault}{WenQuanYi Zen Hei}
% polyglossia no traduce los nombres de listings.
\AtBeginDocument{\providecommand{\lstlistingname}{}\renewcommand{\lstlistingname}{Listado}%
  \providecommand{\lstlistlistingname}{}\renewcommand{\lstlistlistingname}{Índice de listados}}
\usepackage{amsmath, amssymb}
\usepackage{graphicx}
\usepackage[margin=2.5cm]{geometry}
\usepackage[most]{tcolorbox}
\usepackage{etoolbox}
\usepackage{listings}
\usepackage{hyperref}
\usepackage{booktabs}
\usepackage{subcaption}
\usepackage{float}

% Highlight boxes
\newtcolorbox{knowledgebox}[1]{
    enhanced, colback=blue!5!white, colframe=blue!75!black, colbacktitle=blue!75!black,
    coltitle=white, fonttitle=\bfseries, title={#1}, attach boxed title to top left={yshift=-2mm, xshift=2mm},
    boxrule=1pt, sharp corners
}
\newtcolorbox{importantbox}[1]{
    enhanced, colback=yellow!10!white, colframe=yellow!80!black, colbacktitle=yellow!80!black,
    coltitle=black, fonttitle=\bfseries, title={#1}, sharp corners
}
\newtcolorbox{warningbox}[1]{
    enhanced, colback=red!5!white, colframe=red!75!black, colbacktitle=red!75!black,
    coltitle=white, fonttitle=\bfseries, title={#1}, sharp corners
}

\lstset{
    language=Python, basicstyle=\ttfamily\small, keywordstyle=\color{blue},
    stringstyle=\color{red!60!black}, commentstyle=\color{green!60!black},
    breaklines=true, frame=single, numbers=left, numberstyle=\tiny\color{gray},
    captionpos=b, extendedchars=true
}

\newcommand{\notetitle}{Diferencias de ReAct en LangChain/LangGraph y monitoreo con LangSmith}
\newcommand{\noteauthors}{Elaboradas a partir de materiales públicos del curso}
\newcommand{\notedate}{2025}
\newcommand{\videochannel}{Wudaokou Naishi (五道口纳什)}
\newcommand{\videopublishdate}{2025}
\newcommand{\videoduration}{aproximadamente 17 minutos}
\newcommand{\videourl}{https://www.bilibili.com/video/BV1n4n7znEDN}
\newcommand{\videocoverpath}{cover.jpg}
\newcommand{\repourl}{https://github.com/hqhq1025/ai-course-notes}


% Footer with repo info
\usepackage{fancyhdr}
\pagestyle{fancy}
\fancyhf{}
\fancyfoot[C]{\thepage}
\fancyfoot[R]{\footnotesize\color{gray}\href{\repourl}{AI Course Notes}}
\renewcommand{\headrulewidth}{0pt}
\renewcommand{\footrulewidth}{0pt}

\begin{document}

\begin{titlepage}
\centering
{\Large Notas del curso\par}\vspace{1.2cm}
{\huge\bfseries \notetitle\par}\vspace{0.8cm}
{\large \noteauthors\par}\vspace{0.3cm}
{\large \notedate\par}\vspace{1.2cm}
\ifdefempty{\videocoverpath}{}{
\includegraphics[width=0.82\textwidth,height=0.45\textheight,keepaspectratio]{\videocoverpath}\par}
\vfill
\begin{tcolorbox}[width=0.9\textwidth, colback=black!2!white, colframe=black!60, sharp corners]
\textbf{Autor o canal del video}: \videochannel\par
\textbf{Fecha de publicación}: \videopublishdate\par
\textbf{Duración del video}: \videoduration\par
\ifdefempty{\videourl}{}{\textbf{Enlace del video}: \href{\videourl}{\nolinkurl{\videourl}}\par}\par
\textbf{Página del proyecto}: \href{\repourl}{\nolinkurl{\repourl}}\par
\end{tcolorbox}
\end{titlepage}

\tableofcontents
\newpage

\section{Repaso del paradigma ReAct}

\subsection{De CoT a ReAct}

ReAct es uno de los trabajos más emblemáticos en el campo de los agentes (equipo de Yao Shunyu (姚顺雨), publicado a finales de 2022 / principios de 2023 en NeurIPS, con más de 4,000 citas). Su línea de evolución es la siguiente:

\begin{enumerate}
    \item \textbf{CoT (Chain of Thought)}: orientado a responder preguntas: ``Let's think step by step'', que inspiró el Prompt Engineering posterior.
    \item \textbf{Act Only}: el modelo de lenguaje puede interactuar con el entorno, produciendo una acción (Action) y recibiendo del entorno una observación (Observation), pero sin un proceso de pensamiento explícito.
    \item \textbf{ReAct = Reasoning + Acting}: antes de producir la acción se realiza un razonamiento (Thought), lo que reduce la alucinación del modelo y logra un mejor grounding hacia acciones razonables.
\end{enumerate}

\begin{importantbox}{El ciclo central de ReAct}
Thought $\rightarrow$ Action $\rightarrow$ Observation $\rightarrow$ Thought $\rightarrow$ ... $\rightarrow$ Final Answer

ReAct inspiró el diseño y desarrollo de los agentes posteriores, y posiblemente también inspiró la aparición de la característica de Function Calling en la API de GPT (junio de 2023).
\end{importantbox}

\begin{knowledgebox}{Sobre los Reasoning Model}
Los Reasoning Model modernos ya han aprendido a pensar antes de responder una pregunta, sin necesidad de indicarles explícitamente mediante el prompt que «primero deben pensar». Más adelante, al presentar los Reasoning Model, se profundizará en sus normas y estándares de uso.
\end{knowledgebox}

\subsection{Resumen de la sección}
ReAct combina el razonamiento (Reasoning) con la acción (Acting), y constituye la base del diseño de agentes. Entender su tríada Thought-Action-Observation es el requisito previo para entender todos los marcos de agentes posteriores.

\section{Implementación de un ReAct Agent con LangChain}

\subsection{Diseño de la plantilla de Prompt}

En LangChain, \texttt{create\_react\_agent} usa una plantilla de Prompt estandarizada, cuyo diseño enfatiza la \textbf{generalidad y la universalidad}:

\begin{quote}
\textit{Answer the following questions as best you can. You have access to the following tools: \{tools\}. Use the following format: Question / Thought / Action / Action Input / Observation / ... / Final Answer. Begin! Question: \{input\} \{agent\_scratchpad\}}
\end{quote}

La plantilla contiene dos marcadores dinámicos:
\begin{itemize}
    \item \texttt{tools}: la lista de herramientas disponibles
    \item \texttt{agent\_scratchpad}: el historial que se mantiene dinámicamente durante la interacción
\end{itemize}

\subsection{Proceso de construcción en tres pasos}

\begin{enumerate}
    \item Declarar el modelo de lenguaje (por ejemplo, GPT-4.1 Nano)
    \item Llamar a \texttt{create\_react\_agent} para crear el Agent
    \item Colocar el Agent dentro de un \texttt{AgentExecutor} para ejecutarlo
\end{enumerate}

\subsection{Mecanismo de funcionamiento: iteración basada en una plantilla de texto}

La implementación de ReAct en LangChain tiene una característica importante: \textbf{no usa una lista de mensajes de varios turnos (message list)}, sino que va agregando toda la información del historial a \textbf{un mismo mensaje human}:

\begin{warningbox}{Forma en que LangChain ReAct administra el historial}
Cada nuevo turno de interacción es una nueva llamada al LLM. La información del historial (Thought, Action, Observation) se agrega a \texttt{agent\_scratchpad} y se envía al modelo como parte de un único mensaje human. Esto significa que no hay una estructura de mensajes de varios turnos: todo el contexto queda apretado en un solo mensaje.
\end{warningbox}

\subsection{Resumen de la sección}
La implementación de ReAct en LangChain administra el historial de la interacción mediante una plantilla de Prompt más la concatenación de texto; la estructura es simple, pero carece de la organización clara de los mensajes de varios turnos.

\section{Implementación de ReAct Agent en LangGraph}

\subsection{Arquitectura basada en grafos}

La implementación de \texttt{create\_react\_agent} de LangGraph es más moderna. Basta con definir el modelo y las herramientas:

\begin{importantbox}{Lo esencial para entender LangGraph: los bordes virtuales (bordes condicionales)}
El ReAct Agent de LangGraph contiene dos nodos centrales:
\begin{itemize}
    \item \textbf{Nodo Agent}: se encarga de producir las solicitudes de llamada a herramientas
    \item \textbf{Nodo Tools}: se encarga de ejecutar las herramientas y devolver los resultados
\end{itemize}
El nodo Agent tiene un \textbf{borde condicional} que decide si se detiene la llamada a herramientas y se entrega directamente la respuesta, o si se continúa llamando a herramientas.
\end{importantbox}

\subsection{Gestión basada en una lista de mensajes}

A diferencia de LangChain, LangGraph gestiona la interacción por completo mediante \textbf{mensajes de conversación de múltiples turnos}:

Flujo de mensajes: Human Message $\rightarrow$ AI Message (con tool\_calls) $\rightarrow$ Tool Message(s) $\rightarrow$ AI Message (respuesta final)

\begin{knowledgebox}{Llamadas paralelas a herramientas}
LangGraph admite \textbf{llamadas paralelas a herramientas}: el campo \texttt{tool\_calls} de un AI Message es una lista, por lo que se pueden llamar varias herramientas al mismo tiempo. Por ejemplo, al consultar «la ciudad natal del campeón del Abierto de Australia», el modelo podría llamar en paralelo dos veces a la herramienta de búsqueda.
\end{knowledgebox}

\subsection{Resumen de la sección}
La implementación de ReAct en LangGraph se basa en una estructura de grafo y una lista de mensajes de múltiples turnos, con una arquitectura más clara y compatibilidad con llamadas paralelas a herramientas.

\section{Plataforma de monitoreo LangSmith}

\subsection{Configuración y uso}

LangSmith es una plataforma de monitoreo en línea, usada para rastrear y depurar el flujo de trabajo del Agent. Modo de configuración:
\begin{enumerate}
    \item Configurar la API Key de LangSmith en \texttt{.env}
    \item Activar la función de tracing (ponerla en true)
    \item Configurar el nombre del project
\end{enumerate}

\subsection{Método central de depuración}

\begin{importantbox}{Tips centrales para el debug}
Entender su flujo de trabajo revisando las \textbf{entradas y salidas} del modelo. Concéntrate en la entrada y la salida de cada llamada al modelo de lenguaje (la parte naranja); entender la entrada y la salida es entender todo el flujo de trabajo. Todos los pasos previos existen para construir esa lista de mensajes.
\end{importantbox}

\subsection{Resumen de la sección}
LangSmith ofrece una capacidad de depuración visual del Agent; el método central consiste en entender todo el flujo de trabajo monitoreando la entrada y la salida de las llamadas al LLM.

\section{Análisis comparativo entre LangChain y LangGraph}

\begin{table}[H]
\centering
\begin{tabular}{lll}
\toprule
\textbf{Dimensión} & \textbf{LangChain} & \textbf{LangGraph} \\
\midrule
Gestión del historial & Texto añadido a un solo mensaje & Lista de mensajes de varios turnos \\
Estructura de mensajes & Un solo mensaje human (con scratchpad) & Secuencia de mensajes Human / AI / Tool \\
Llamadas a herramientas & En serie (una a la vez) & Con soporte para paralelismo (lista de tool\_calls) \\
Arquitectura & Plantilla de prompt + Executor & Graph (nodos + bordes condicionales) \\
Visualización & Requiere seguimiento manual & Se puede exportar directamente el diagrama del Graph \\
\bottomrule
\end{tabular}
\caption{Comparación de la implementación de ReAct en LangChain y LangGraph}
\end{table}

\subsection{Resumen de la sección}
LangGraph supera de manera integral a la implementación anterior de LangChain en claridad de arquitectura, gestión de mensajes y capacidad de paralelismo. El enfoque de LangChain resulta más adecuado para entender el concepto original de ReAct, mientras que LangGraph es más apropiado para entornos de producción.

\section{Síntesis y ampliación}

\begin{enumerate}
    \item \textbf{Dar importancia a los conceptos básicos}: entre más básico es algo, más esencial resulta; todas las ideas modernas y potentes de desarrollo de Agentes que vinieron después nacieron de los trabajos clásicos.
    \item \textbf{Entender las diferencias de implementación}: LangChain funciona mediante gestión de plantillas y concatenación de texto, mientras que LangGraph funciona mediante mensajes de múltiples turnos y una estructura de grafo; los mecanismos subyacentes son distintos, y entender la diferencia ayuda a elegir la herramienta adecuada.
    \item \textbf{Aprovechar bien las herramientas de monitoreo}: plataformas como LangSmith son un recurso valioso para entender el flujo de trabajo interno de un Agent; lo esencial es revisar la entrada y la salida de cada llamada al LLM.
\end{enumerate}

\end{document}
